\section*{Capítulo \ref{ch:b2:dipole-radiation} --- Radiación dipolar y dispersión}

\begin{solution}{exo:b2:dipole-radiation:1}
$p_0 = I_0\ell/\omega = \qty{3.2e-9}{C.m}$; $E = \mu_0\omega^2p_0\sin\theta/4\pi r$: \qty{0.13}{V/m}
en $\theta = 90^\circ$ y \qty{0.063}{V/m} a \qty{30}{\degree}; $\mathcal P = \mu_0\omega^4p_0^2/12\pi c =
\qty{175}{W}$; $R_{\text{rad}} = 2\mathcal P/I_0^2 = \qty{88}{\ohm}$.
\end{solution}

\begin{solution}{exo:b2:dipole-radiation:2}
$790(\ell/\lambda)^2$: \qty{8.8e-5}{\ohm}, \qty{8.8e-3}{\ohm}, \qty{0.88}{\ohm}, \qty{88}{\ohm}. Cobre:
$0.026$, $0.084$, $0.27$, \qty{0.84}{\ohm} (\omterm{thm:b2:waves-in-media:skin}{efecto pelicular}): rendimientos del $0.3\%$, $10\%$,
$77\%$ y $99\%$. En longitudes de onda largas, solo un mástil que sea una fracción apreciable de
$\lambda$ tiene una \omterm{ex:b2:dipole-radiation:antenna}{resistencia de radiación} superior a sus pérdidas.
\end{solution}

\begin{solution}{exo:b2:dipole-radiation:3}
$400 : 550 : 700 = 9.4 : 2.6 : 1$. $(10/3)^4 = 120$: el radar de \qty{3}{cm}
ve las gotas pequeñas cien veces mejor; el de \qty{10}{cm} mira
a través de la lluvia a lo que hay detrás.
\end{solution}

\begin{solution}{exo:b2:dipole-radiation:4}
$\sin^2\chi/(1 + \cos^2\chi)$: $0.14$, $0.60$, $1$, $0.60$. Apúntese a \qty{90}{\degree} del
Sol (un círculo máximo por el cielo); con el Sol en el cenit, todo el
horizonte está a \qty{90}{\degree} y polarizado a lo largo de él. La dispersión múltiple,
el reflejo del suelo y la anisotropía de las moléculas añaden luz sin
polarizar: alrededor del $75\%$ como mucho.
\end{solution}

\begin{solution}{exo:b2:dipole-radiation:5}
(a) $\int\sin^2\theta\,\dd\Omega = 8\pi/3$. (b) $\int_{60^\circ}^{120^\circ}\sin^3\theta\,\dd\theta/
(4/3) = 0.917/1.333 = 69\%$. (c) $\sin^2\theta = \tfrac12$: $\theta = \qty{45}{\degree}$, un
haz de \qty{90}{\degree} de anchura. (d) Máximo $1$ frente a la media $2/3$: $1.5$.
\end{solution}

\begin{solution}{exo:b2:dipole-radiation:6}
(a) $p/4\pi\varepsilon_0r^3 = \omega^2p/4\pi\varepsilon_0c^2r$ en $r_0 = c/\omega = \lambda/2\pi$. (b) \qty{48}{m}:
a \qty{10}{m}, el campo cuasiestático; a \qty{10}{km}, el radiado. (c)
$r_0 = \qty{4.8}{cm}$: la cabeza está en la zona próxima. (d) En la zona próxima, $\vect E$
y $\vect B$ están en cuadratura: el \omterm{thm:b2:poynting-vector:poynting}{vector de Poynting} promedia cero,
la energía se almacena y se devuelve, como en un condensador.
\end{solution}

\begin{solution}{exo:b2:dipole-radiation:7}
(a) $8\pi r_e^2/3 = \qty{6.65e-29}{m^2}$. (b) $n_e\sigma_TL = 10^{14} \times 6.65 \times 10^{-29}
\times 10^9 = 7 \times 10^{-6}$: una millonésima de la fotosfera, sin color porque
la \omterm{prop:b2:dipole-radiation:rayleigh}{dispersión Thomson} tampoco lo tiene, y ahogada por la luz solar dispersada del cielo
salvo en un eclipse. (c) \qty{10}{keV} $\gg$ \qty{300}{eV}: los electrones responden como
libres, $6\sigma_T = \qty{4e-28}{m^2}$. (d) La absorción fotoeléctrica crece como $Z^4$
y destaca el calcio; la dispersión es la misma por electrón en el hueso
y en la carne.
\end{solution}

\begin{solution}{exo:b2:dipole-radiation:8}
(a) $\tau = 0.36$, $0.10$, $0.038$: $T = 0.70$, $0.90$, $0.96$. (b) $\tau \times 38 = 13.6$,
$3.8$, $1.45$: $T = 1 \times 10^{-6}$, $0.02$, $0.23$; rojo entre azul $\sim 2 \times 10^5$. (c)
La atmósfera terrestre curva hacia la sombra la luz de la puesta de sol, enrojecida
por esa misma dispersión. (d) Cerca del horizonte el camino es largo: el
azul se dispersa fuera de la propia luz dispersada, y la dispersión
múltiple añade blanco.
\end{solution}

\begin{solution}{exo:b2:dipole-radiation:9}
(a) A lo largo de la línea del par, la diferencia de camino $\lambda/2$ da
cancelación; transversalmente, suma. (b) En antifase, al revés:
radiación según la línea (longitudinal) y nada transversalmente. (c) Hacia $+x$:
retardo de camino $-\pi/2$ más la fase de alimentación $+\pi/2$, en fase; hacia $-x$:
$-\pi$, que cancela: una cardioide, de un solo lado. (d) Los mástiles conforman la cobertura
y protegen a los vecinos; una agrupación en fase orienta el haz cambiando las fases
de alimentación, en microsegundos.
\end{solution}

\begin{solution}{exo:b2:dipole-radiation:10}
(a) $\langle\mathcal P\rangle = e^2\omega_0^4x_0^2/12\pi\varepsilon_0c^3$, $W = \tfrac12m\omega_0^2x_0^2$:
$\Gamma = \mathcal P/W = e^2\omega_0^2/6\pi\varepsilon_0mc^3$. (b) $\omega_0 = \qty{3.2e15}{rad/s}$: $\Gamma =
\qty{6.4e7}{s^{-1}}$ --- el valor usado. (c) \qty{16}{ns}; $\Delta\lambda = \lambda^2\Gamma/2\pi c =
\qty{1.2e-14}{m}$, una centésima de picómetro (\qty{10}{MHz}). (d) $Q = 5 \times 10^7$.
\end{solution}

\begin{solution}{exo:b2:dipole-radiation:11}
(a) \qty{75}{m}; $I_0 = \sqrt{2\mathcal P/R} = \qty{53}{A}$. (b) $\langle\Pi\rangle = 1.5 \times
5 \times 10^4/2\pi \times 10^8 = \qty{1.2e-4}{W/m^2}$, $E_0 = \sqrt{2\Pi/\varepsilon_0c} = \qty{0.30}{V/m}$.
(c) \qty{0.30}{V}. (d) Las corrientes de retorno circulan por el suelo; un
terreno resistivo añade pérdidas y estropea el espejo; los radiales enterrados les dan un
camino de cobre.
\end{solution}

\begin{solution}{exo:b2:dipole-radiation:12}
(a) $p_0 = \alpha E_0$ da $\sigma = \alpha^2\omega^4/6\pi\varepsilon_0^2c^4$; con $\alpha = \varepsilon_0(n^2
- 1)/N$ y $\omega = 2\pi c/\lambda$: $\sigma = 8\pi^3(n^2 - 1)^2/3N^2\lambda^4$. (b) $n^2 - 1 =
5.8 \times 10^{-4}$: $\sigma = \qty{4.9e-31}{m^2}$, recorrido libre medio $1/N\sigma = \qty{80}{km}$. (c)
$8/80 = 0.1$. (d) $\sigma$ se mide por la \omterm{def:b2:dispersion-wave-packets:complex}{atenuación} del cielo y $n$ es
conocido: se deduce $N$ y, con él, el número de Avogadro.
\end{solution}

\begin{solution}{pb:b2:dipole-radiation:1}
\textbf{1.} $\lambda = \qty{300}{m}$, $h = \qty{75}{m}$; $I_0 = \sqrt{2 \times 10^5/36} = \qty{75}{A}$;
$V = RI_0 = \qty{2.7}{kV}$.

\textbf{2.} $\ell_{\text{ef}} = \qty{48}{m}$: $p_0 = 75 \times 48/6.3 \times 10^6 = \qty{5.7e-4}{C.m}$.

\textbf{3.} $\langle\Pi\rangle = 1.5 \times 10^5/2\pi r^2$: \qty{2.4e-4}{W/m^2} y \qty{2.4e-6}{W/m^2};
$E_0 = \qty{0.42}{V/m}$ y \qty{0.042}{V/m}.

\textbf{4.} Látigo: \qty{42}{mV}. Varilla: $N\mu_{\text{ef}}A\omega B_0 = 10^4 \times 10^{-4} \times 6.3 \times
10^6 \times 1.4 \times 10^{-10} = \qty{0.9}{mV}$ --- menos, pero compacta y selectiva.

\textbf{5.} $A = 1.5\lambda^2/4\pi = \qty{1.1e4}{m^2}$: $P = \Pi A = \qty{26}{mW}$.

\textbf{6.} $36/38 = 95\%$; \qty{5}{kW} de calor.

\textbf{7.} $r/\lambda = 33$ a \qty{10}{km}: \omterm{thm:b2:dipole-radiation:field}{zona de radiación}; a \qty{100}{m},
$r \approx \lambda/3$: la región de transición, donde el campo cuasiestático todavía
cuenta.

\textbf{8.} De día la capa D absorbe la onda que si no llegaría
a las capas reflectantes; de noche desaparece y la onda ionosférica
lleva la emisora a cientos de kilómetros.

\textbf{9.} $\sigma(450) = \qty{1.0e-30}{m^2}$, $\sigma(650) = \qty{2.4e-31}{m^2}$: cociente $4.4$.

\textbf{10.} Recorridos libres medios de $39$, $87$ y \qty{170}{km}; $\tau = 0.21$, $0.09$, $0.05$:
dispersado el $19\%$, el $9\%$ y el $5\%$.

\textbf{11.} $N\sigma I = 2.5 \times 10^{25} \times 4.6 \times 10^{-31} \times 300 = \qty{3.5}{mW/m^3}$,
en $4\pi$ con el diagrama $(1 + \cos^2\chi)$ para \omterm{def:b2:plane-waves-polarization:states}{luz natural}.

\textbf{12.} La columna dispersa $3.5 \times 10^{-3} \times 8000 = \qty{28}{W}$ de sus
\qty{300}{W} (el $9\%$); repartidos por el hemisferio, el cielo entrega unos
\qty{4}{W.m^{-2}.sr^{-1}}, frente a los $300/6.8 \times 10^{-5} = \qty{4e6}{W.m^{-2}.sr^{-1}}$ del Sol:
una millonésima por estereorradián --- el orden correcto.

\textbf{13.} Totalmente polarizada a \qty{90}{\degree} en dispersión simple; la dispersión
múltiple y el suelo la reducen a un $75\%$; las abejas y las hormigas leen
el patrón para encontrar el Sol tras las nubes.

\textbf{14.} $\tau \times 38$: azul $8$ ($T = 3 \times 10^{-4}$), rojo $1.8$ ($T = 0.17$): un Sol
rojo intenso; en el cenit la luz llega tras un camino corto y es azul.

\textbf{15.} Sin atmósfera no hay nada que disperse: el cielo es negro junto
al Sol.

\textbf{16.} El polvo de un micrómetro dispersa como partículas de Mie: hacia delante y
en rojo, lo que da un cielo de color caramelo; cerca del Sol poniente, la
luz dispersada hacia delante es azul.

\textbf{17.} $8\pi r_e^2/3 = \qty{6.65e-29}{m^2}$.

\textbf{18.} $a = v^2/2d = \qty{6.4e21}{m/s^2}$; $\mathcal P = e^2a^2/6\pi\varepsilon_0c^3 =
\qty{2e-10}{W}$ durante $2d/v = \qty{2.5e-14}{s}$: \qty{6e-24}{J}, $10^{-10}$ de
\qty{100}{keV} --- el rendimiento real ronda el $1\%$, porque el frenado ocurre
en encuentros violentos y próximos con los núcleos, no en una deceleración
uniforme y suave.

\textbf{19.} $hc/E = \qty{12.4}{pm}$.

\textbf{20.} \qty{2}{kW} en el haz y unos \qty{20}{W} de rayos X; dos kilovatios
de calor en un punto --- el ánodo gira para repartirlo.

\textbf{21.} $a = v^2/r = \qty{9e22}{m/s^2}$, $\mathcal P = \qty{5e-8}{W}$; los \qty{13.6}{eV} se van
en \qty{5e-11}{s}.

\textbf{22.} Las cargas clásicas en movimiento ligado radian sin parar;
la mecánica cuántica tiene estados estacionarios que no lo hacen, y uno más bajo
sin sitio al que caer.

\textbf{23.} $a = c^2/R = \qty{9e14}{m/s^2}$: \qty{5e-24}{W} por Larmor, y $\gamma^4 \sim 10^{16}$
veces más en la realidad: decenas de kilovatios de luz de sincrotrón de un
haz almacenado.

\textbf{24.} $\sigma_TI = \qty{6.7e-26}{W}$; $1.5 \times 10^{25}$ electrones para un vatio.

\textbf{25.} Mástil: \qty{5.7e-4}{C.m} a \qty{1}{MHz}, \qty{100}{kW}, un toro
vertical. Molécula: $\sim$\qty{1e-35}{C.m} a \qty{500}{THz}, \qty{1e-28}{W},
un toro en torno al campo. Electrón libre: Thomson, \qty{1e-25}{W} bajo
la luz solar, con el mismo diagrama.
\end{solution}
